\section{Discussion}

Finite disturbances in this spatial feedback game are governed by two related but distinct aspects of spatial organization. First, the initial arrangement of defectors determines the cooperator--defector contacts available for strategy updates. Dispersed seeds have a longer, initially expansion-biased interface and are more likely to reach the early growth threshold than compact seeds with the same size and resource deficit. Second, as a disturbance grows, the evolving strategies create a depleted interior and an interfacial resource layer. The position of that layer relative to the strategy boundary affects the probability of subsequent collapse. Resource translations separate this alignment effect from changes in strategy geometry or in the resource field's own distribution and autocorrelation.

This second result identifies a limitation of descriptions based on separate statistics of the two fields. Network reciprocity explains how the arrangement of strategic interactions affects selection \cite{szabo2007graphs,perc2017cooperation}, while local environmental feedback makes the payoffs themselves depend on resources generated or depleted by the competing populations \cite{lin2019spatial,chen2025localized}. In the present model, states with identical strategy configurations and identical single-field resource statistics can nevertheless have different finite-time collapse probabilities. The strategy--resource covariance and the distribution of local interfacial payoff gaps provide complementary descriptions of this relative alignment, although neither has been shown to determine collapse probability on its own.

The exact relations clarify different parts of the mechanism. The first-visit identity separates the frequency of strategy-changing events from their directional bias. The resource-balance identity shows how strategy--resource covariance contributes to the next mean-resource increment. The sensitivity calculation identifies where a resource perturbation most strongly affects the signed interface bias. These statements are exact for their specified updates and assumptions; the measured collapse probabilities reflect subsequent stochastic evolution. In particular, reversing the first-step bias by rearranging boundary resources does not guarantee a corresponding reversal in later population trajectories or in the finite-time collapse probability.

These observations connect the game to the physics of interfaces evolving through a self-modified medium. Theories of front propagation describe how local rules govern moving boundaries \cite{vansaarloos2003fronts}, and spatial ecological models show how environmental feedback interacts with domain growth and collapse \cite{bel2012gradual,rietkerk2021tipping}. Here the resource field relaxes toward a configuration-dependent target, but strategy changes continually shift that target. The breakdown of a simple ordering between initial interface response and later dynamics shows why a description based only on instantaneous bias is incomplete. A reduced interfacial theory would need to retain the relevant environmental response or justify its elimination through a controlled approximation.

The scope of these conclusions is finite-time and model-specific. The formation of a dominant defector cluster does not establish a universal coarsening mechanism; distinguishing curvature-driven motion from other sources of interface evolution would require additional length-scale and size-dependent measurements \cite{bray1994ordering,dornic2001coarsening}. Likewise, the observed collapse crossovers are defined by first-passage criteria rather than by thermodynamic phase boundaries. The perturbations are externally imposed on an absorbing cooperative background, so the present experiments do not identify a spontaneous nucleation rate or a critical droplet of the kind considered in metastable decay and stochastic escape theory \cite{langer1969decay,assaf2017wkb}. Some resource-alignment contrasts also weaken or become statistically unresolved under changes in system size or depletion threshold, as reported in Supplemental Material, Sec.~\ref{sm-sec:checks}.

The principal finding is therefore not a universal collapse threshold, but a controlled separation of two contributors to finite-disturbance response: the geometry of the initial strategy interface and its alignment with a dynamically generated resource field. Further comparisons at matched interface activity and initial bias, together with systematic variation of disturbance and system sizes, could test which aspects of this coupled spatial history persist beyond the parameter regime studied here.
